\makeatletter
\def\input@path{{../styles/}{../../styles/}{../../../styles/}{../}{../../}{../../../}}
\makeatother


\documentclass{ee102_pset}
\input{macros.tex}
\input{preamble.tex}
\usepackage{amsmath}
\usepackage{amssymb}
\usepackage{graphicx}
\usepackage{listings}
\usepackage[american]{circuitikz}

% Assignment info
\author{\rule{3cm}{0.4pt}} % Name placeholder
\submitdate{\rule{3cm}{0.4pt}} % Submission date placeholder
\problemset{Problem Set \#5: Applying Convolution}
\renewcommand{\duedate}{October 7, 2026}
\shorttitle{Problem Set \#5}

\begin{document}
\problem{1}
A system responds to an impulse input $\delta(t)$ in an exponentially decaying manner. So, the impulse response of the system is given by:
\[
h(t) = e^{-2t}u(t)
\]
where $u(t)$ is the unit step function.
\problempart [10 points] Is the system linear and time-invariant? Justify your answer.
\problempart[10 points] What is the output $y(t)$ of the system when the input is $x(t) = k \delta(t)$, where $k$ is a constant?

\problempart[20 points] What is the output $y(t)$ of the system when the input is $x(t) = u(t)$? You must find this by starting from the relationship between the step signal and the impulse signal (from problem set \#2): 
\[
u(t) = \int_{-\infty}^{t} \delta(\tau) d\tau
\]

Then, use linearity and the impulse response $h(t)$ to find the output $y(t)$.

\problempart[20 points] Find the output $y(t)$ of the system when the input is a pulse signal of amplitude $A$ and duration $\tau, \quad x(t) = A[u(t) - u(t-\tau)]$.

\problempart[10 points each] Find the output of the system to the following input signals
\begin{enumerate}
    \item $x(t) = e^{-3t}u(t)$
    \item $x(t) = \cos(2t)u(t)$
    \item $x(t) = \sin(4t)u(t)$
    \item $x(t) = t u(t)$
\end{enumerate}


\end{document}